\part{Paquete de plantilla.}

\chapter{Modulo de Secciones.}

\section{Resaltado de texto.}

    \subsection{Término nuevo.}
    
    \begin{description}
        \item[¿Qué se considera un término nuevo?.]\hfil
        \begin{itemize}
            \item Conceptos clave. Palabras propias del tema que el autor está explicando o definiendo (ej. taxonomía, hipervínculo).
            \item Tecnicismos o jerga. Palabras especializadas de una disciplina que el lector general puede desconocer.
            \item Neologismos: Palabras de reciente creación en el idioma o adaptaciones recientes.
            \item Extranjerismo. Palabras de otros idiomas.
        \end{itemize}
        \item[Reglas de uso.] \hfil
        \begin{itemize}
            \item Primera aparición:  link al glosario.
            \begin{itemize}
                \item Si no es su definición, el glosario no hace referencia a esta primera aparición.
            \end{itemize}
            \item Los enlaces se resaltan en azul.      
            \item El glosario hace referencia a su definición.    
            \item En la definición, se utiliza un link que va al glosario y se resalta en negritas.
            \item Las siguientes apariciones no se resaltan.              
            \item \textbf{No resaltar con vínculo al glosario hasta después de la definición.} Esto para asegurar que solamente se haga referencia a términos que ya cuentan con una entrada en el glosario completa (con la referencia a su definición).
        \end{itemize}        
        \item[Ejemplos de términos nuevos.]\hfil
        \begin{itemize}
            \item Primera aparición que no es definición. \glshyperlink{API}.
            \item Definición: siga el link del punto anterior y lo llevará a la definición, que es una entrada de glosario, en negritas.
            \item Después de la primera aparición y aparte de la definición se usan letras comunes.
            \item Los acrónimos, se explica su significado en la definición en el texto y en el glosario.            
        \end{itemize}
    \end{description}
    
    \subsection{Función metalingüistica.}
    % REFERENCIAS
    % https://concepto.de/funcion-metalinguistica/
    % https://www.rae.es/libro-estilo-lengua-espa%C3%B1ola/clase-de-letra#13
    % https://marianruiz.com/cursiva-o-comillas/
    % https://www.fundeu.es/wp-content/uploads/2013/05/CursivasGuiaFundeu.pdf
    
    
    La función metalingüistica del lenguaje se da cuando utilizamos la lengua para hablar del propio lenguaje. Esta función se centra en el código, es decir, en la lengua y sirve para:
    
    \begin{itemize}
        \item Analizar el lenguaje. Explicar el funcionamiento de una lengua desde alguna disciplina lingüistica, como la ortografía, la sintaxis o la semántica.
        \item Realizar definiciones. Se usa para definir palabras.
        \item Se usa para traducir términos, expresiones u oraciones en otra lengua.
        \item Explicar algo que se mencionó previamente. Se usa para aclarar algo que se dijo antes.
    \end{itemize}
    
        \subsubsection{Reglas.}
        
        Clave para distinguir el uso normal y uso metalingüistico.
        
        \begin{itemize}
            \item Preguntarte: ¿puedo reemplazar la palabra por su definición y la oración sigue teniendo sentido?
            \begin{itemize}
                \item \textquote{El \textbf{sprint} dura dos semanas} - \textquote{El \textbf{periodo corto de trabajo} dura dos semanas}: Tiene sentido, es uso.
                \item \textquote{El término \textbf{sprint} proviene del inglés} - \textquote{El término \textbf{periodo corto de trabajo} proviene del inglés}: No tiene sentido, es mención.
            \end{itemize}
            \item Si puedes señalar el término y decir \textquote{esta secuencia de letras...} o \textquote{esta palabra que se escribe así...}, es mención y va en cursiva. Si el término está funcionando normalmente dentro de la oración (refiriéndose a su significado), es uso y va en redondas.
        \end{itemize}
        
        \subsubsection{Ejemplos.}
    
   
        
        \begin{itemize}
            \item \textit{Abrir} es el antónimo de \textit{cerrar}, por que sus significados son opuestos.
            \item ¿Qué significa la palabra \textit{átomo}?
            \item Cuando dije \textit{naranja}, me refería a la fruta, no al color.
            \item En inglés, \textit{vaca} se dice \textit{cow}.
            \item \textit{Cocinar} es un infinitivo.
            \item Un adjetivo es una palabra que designa una característica o cualidad.
            \item ¿Cómo se pronuncia \textit{hat} en inglés?
            \item \textit{Metáfora} es una palabra esdrújula por que su acento recae en la antepenúltima sílaba.
            \item \textit{Hardware} es un anglicismo, es decir, un término que proviene del inglés.
            \item \textit{México} es un sustantivo propio.
        \end{itemize}
        
        \begin{itemize}
            % "sprint" es extranjerismo + mención → cursiva por ambas razones
            \item El término \textit{sprint} proviene del inglés.
            
            % "API" no es extranjerismo, pero aquí hablas sobre la sigla como objeto lingüístico → cursiva
            % "Application Programming Interface" es mención de cómo se dice en inglés → cursiva
            \item La sigla \textit{API} significa \textit{Application Programming Interface}.
            
            % "branch" y "fork" son extranjerismos + mención → cursiva
            \item No confundir \textit{branch} con \textit{fork}.
            
            
            % === USO NORMAL (usas la palabra con su significado) ===
            
            % "sprint" es extranjerismo → cursiva (no por mención, sino por extranjerismo)
            \item El \textit{sprint} dura dos semanas.
            
            % "API" no es extranjerismo → redondas, la estás usando normalmente
            \item Se conecta a la API del servicio.
            
            % "branch" es extranjerismo → cursiva (por extranjerismo, no por mención)
            \item Creó un \textit{branch} para la nueva funcionalidad.
            
            % La palabra "murciélago" no se refiere al animal, sino a la palabra como objeto.
            % Prueba: "La palabra [animal volador nocturno] contiene todas las vocales" → no tiene sentido.
            % → Es mención → cursiva.
            \item La palabra \textit{murciélago} contiene todas las vocales.
            
            % Aquí preguntas por el significado de la palabra, no la estás usando para referirte
            % a la capacidad de adaptarse. Hablas sobre la palabra en sí.
            % Prueba: "¿Qué significa [capacidad de adaptarse]?" → no tiene sentido.
            % → Es mención → cursiva.
            \item ¿Qué significa la palabra \textit{resiliencia}?
            
            % "esdrújulas" aquí se refiere a una categoría de palabras, no a una palabra específica
            % que estés mencionando como objeto. Estás USANDO el término con su significado normal.
            % Prueba: "Las [palabras con acento en la antepenúltima sílaba] siempre llevan tilde" → sí tiene sentido.
            % → Es uso → redondas.
            \item Las palabras esdrújulas siempre llevan tilde.
            
            % "-ar" no es una palabra que estés usando con significado propio, es un fragmento
            % lingüístico que estás mostrando como objeto. No puedes sustituirlo por su definición.
            % → Es mención → cursiva.
            \item En español, los verbos terminados en \textit{-ar} pertenecen a la primera conjugación.
            
        \end{itemize}
    
    \subsection{Convenciones.}
    
    \begin{description}
        \item[Término nuevo.]\hfil
        \begin{itemize}
            \item Agregar enlace al glosario resaltado en azul y negritas al definirlo.
            \item Si la primera mención del término fué antes de su definición. Solo cuando ya existe la definición y la entrada del glosario, agregar una referencia al glosario en la primer mención. La entrada del glosario no deberá hacer referencia a la primera mención, sino mantener la referencia que ya tenía a la definición.
            \item Los términos de uso común en el texto y asentados en la disciplina del texto, aunque sean extranjerismos, no se resaltarán con cursivas.
            \item Términos que no son de uso común en el texto o frases completas en otros idiomas se resaltarán con cursivas.
        \end{itemize}
    \end{description}

\begin{description}
    \item[Término nuevo (primera vez)].
\end{description}

\begin{itemize}
    \item Extranjerismos: \nkForeignism{branch},% \fr{commit}.
    \item Nombres de herramientas: Git, SonarQube.
    \item Términos nuevos (primera vez): Conceptos clave,  azul + link al glosario.    

\end{itemize}

El paquete \nkLatexPackage{tcolorbox} permite crear cajas con estilo.



\section{Secciones.}

    \subsection{Objetivo.}

    \begin{nkObjective}
        Comprender las propiedades de los espacios vectoriales.
    \end{nkObjective}
    
    \subsection{Definición.}
    
    \begin{definition}
        \label{def:definition_api}
        Una \textbf{\gls{API}} (\textit{Application Programming Interface} / Interfaz de Pogramación de Aplicaciones) es una pieza de código que permite a dos aplicaciones comunicarse entre sí para compartir información y funcionalidades. Se usan generalmente en bibliotecas de programación.
    \end{definition}
    
    
    
    \begin{definition}[modeling language]
        \label{def:definition_example}
        Los \glslink{modeling language}{machine} se desarrollaron precisamente para estos escenarios y demuestran reglas claramente definidas para descripción estructurada de un sistema. Permiten representar diferentes vistas de la arquitectura.
    \end{definition}
1
\section{Tablas.}

\begin{definition}[Domain-Driven Design]
	\label{def:ddd}
	Domain-Driven Design is an approach to software development that centers the development on programming a domain model that has a rich understanding of the process and rules of a domain.
\end{definition}

\begin{table}[H]
	\centering
	\begin{tblr}{
			width = 0.95\textwidth,
			colspec = {|X[l,m,2]|X[c,m, 1]|X[l,m,4]|}, % Definimos la alineación en la especificación de columnas
			hlines,
			vlines,		
			rows = {m},
			row{odd} = {bg=gray!30},
			%row{even} = {bg=red!30},
			row{1} = {font=\bfseries, bg=gray!60, c}, % La primera fila centrada
		}	
		\textbf{Nombre} & \textbf{Acrónimo} & \textbf{Significado} \\ \hline 
		Java Development Kit & JDK & El software para programadores que desean escribir programas en Java. \\  
		Java Runtime Environment & JRE & El software para consumidores que desean ejecutar programas Java.\\ 		
	\end{tblr}
	\caption{\label{tab:template_table} Ejemplo de tabla formateada.}
\end{table}


\ref{tab:table_example}

\begin{nkTable}{|X[l,m,2]|X[c,m,1]|X[l,m,4]|}{Ejemplo de tabla.}{tab:table_example}
	\textbf{Nombre} & \textbf{Acrónimo} & \textbf{Significado} \\ 
	Java Development Kit & JDK & Software para \linebreak programadores. \\  
	Java Runtime Environment & JRE & Software para ejecutar programas Java. \\    
\end{nkTable}

\begin{nkTable}[0.8\textwidth]{|X[l,m,2]|X[c,m,1]|X[l,m,4]|}{Ejemplo de tabla.}{tab:table_example2}
    \textbf{Nombre} & \textbf{Acrónimo} & \textbf{Significado} \\ 
    Java Development Kit & JDK & Software para programadores. \\  
    Java Runtime Environment & JRE & Software para ejecutar programas Java. \\    
\end{nkTable}




\section{Preguntas}

Ejemplos de uso

Opción múltiple con una respuesta correcta:

\begin{nkQuestion}{¿Qué comando actualiza un submódulo?}
    \nkAnswer[correct]{git submodule update --remote}
    \nkAnswer{git pull --submodules}
    \nkAnswer{git fetch submodule}
\end{nkQuestion}

Opción múltiple con varias correctas:

\begin{nkQuestion}{¿Cuáles son principios SOLID? (selecciona todos los correctos)}
    \nkAnswer[correct]{Single Responsibility}
    \nkAnswer[correct]{Open/Closed}
    \nkAnswer{Don't Repeat Yourself}
    \nkAnswer[correct]{Liskov Substitution}
    \nkAnswer{Keep It Simple}
\end{nkQuestion}

Respuesta única (tipo flashcard):

\begin{nkQuestion}{¿Qué archivo almacena la configuración de submódulos en git?}
    \nkAnswer[correct]{.gitmodules}
\end{nkQuestion}


% Opción múltiple — en modo estudio muestra las opciones sin check/cross
\begin{nkQuestion}{¿Qué comando actualiza un submódulo?}
    \nkAnswer[correct]{git submodule update --remote}
    \nkAnswer{git pull --submodules}
    \nkAnswer{git fetch submodule}
\end{nkQuestion}

% Pregunta abierta — en modo estudio muestra solo la línea en blanco
\begin{nkQuestion}[open]{¿Qué archivo almacena la configuración de submódulos?}
    \nkAnswer[correct]{.gitmodules}
\end{nkQuestion}

    \subsection{Abreviación}
    
    \begin{nkQuestion}[open]{¿Qué archivo almacena la configuración de submódulos?}
        \ans[correct]{.gitmodules}
    \end{nkQuestion}
    
    
\section{Referencias}


\begin{itemize}
    \nkRefcap{IEEEreferencias:Ref1}{Libro 1}
    \nkRefcap{onlineAwsConcepts}{Curso}
\end{itemize}

